\begin{frame}{자주 하는 실수: 풀링을 ``특징 추출''로 오해}
  \begin{block}{정리}
    풀링은 \textbf{학습 파라미터가 없는 순수 다운샘플링} 연산이다.
    ``특징을 추출한다''는 표현은 \textbf{합성곱 층의 역할},
    풀링의 역할은 \textbf{이미 학습된} 특징지도를
    (a) 공간 해상도를 줄이고, (b) 미세한 위치 변동에
    둔감하게 만드는 것.
  \end{block}
  \vskip0.6em
  이 구분을 섞으면 두 가지 흔한 설계 실수:
  \begin{itemize}
    \item \kb{풀링을 너무 일찍, 너무 많이}: 첫 합성곱 직후
      2$\times$2 풀링을 하면 아직 ``선분''을 감지하는 단계인데
      해상도를 1/4로 줄여 \textbf{세부 정보를 잃음}. 실무에서는
      \textbf{첫 두$\sim$세 합성곱 층은 풀링 없이}(스트라이드 1) 쌓고,
      채널이 충분히 늘어난 뒤 풀링
    \item \kb{풀링을 ``합성곱의 대안''으로}: 풀링은
      \textbf{합성곱 이후}에 쓰는 연산이지, \textbf{대신}이 아니다.
      ``3$\times$3 합성곱이 비싸니까 2$\times$2 풀링으로 대체''는
      학습하는 가중치를 \textbf{아예 없애버리는 것}
      (스트라이드 2 합성곱은 \textbf{학습 가능한} 다운샘플링이라
      혼용되기도 하지만, 근본적 차이가 ``학습 여부'')
  \end{itemize}
\end{frame}
